% Part: normal-modal-logic
% Chapter: sequent-calculus
% Section: rules-for-K

\documentclass[../../../include/open-logic-section]{subfiles}

\begin{document}

\olfileid{nml}{seq}{rul}

\olsection{\Log{K} साठी नियम}

नेहमीच्या विधानीय संयोजकांसाठीचे नियम हे
नेहमीच्या क्रमवर्ती कलनातील~$\Log{LK}$ नियमांसारखेच
आहेत. स्वयंसिद्धकेही तीच आहेत: $!A \Sequent !A$
या रूपातील कोणतीही क्रमवर्ती स्वयंसिद्धक असते.

मोडल संकारक\iftag{prvBox}{\iftag{prvDiamond}{~$\Box$
आणि}{~$\Box$}}{}\iftag{prvDiamond}{~$\Diamond$}{} यांच्यासाठी
पुढील अतिरिक्त \iftag{notprvBox,notprvDiamond}{नियम}{नियम}
आहेत:
\[
  \iftag{prvBox}{\iftag{prvDiamond}{% <> and [] primitive
    \Axiom$\Gamma \fCenter \Delta, !A$
    \RightLabel{$\Box$}
    \UnaryInf$\Box\Gamma \fCenter \Diamond\Delta, \Box!A$
    \DisplayProof
    \qquad
    \Axiom$!A, \Gamma \fCenter \Delta$
    \RightLabel{$\Diamond$}
    \UnaryInf$\Diamond!A, \Box\Gamma \fCenter \Diamond\Delta$
    \DisplayProof
}{% only Box primitive
\Axiom$\Gamma \fCenter !A$
\RightLabel{$\Box$}
\UnaryInf$\Box\Gamma \fCenter \Box!A$
\DisplayProof
}}{% only <> primitive
  \Axiom$!A \fCenter \Delta$
  \RightLabel{$\Diamond$}
  \UnaryInf$\Diamond!A \fCenter \Diamond\Delta$
  \DisplayProof
}
\]
येथे \iftag{prvBox}{$\Box\Gamma$ म्हणजे~$\Gamma$
मधील प्रत्येक !!{formula} च्या आधी~$\Box$
लावून~$\Gamma$ पासून मिळणारा !!{formula}s चा अनुक्रम\iftag{prvDiamond}{;
आणि }{}}{}\iftag{prvDiamond}{$\Diamond\Delta$ म्हणजे~$\Delta$
मधील प्रत्येक !!{formula} च्या आधी~$\Diamond$
लावून~$\Delta$ पासून मिळणारा !!{formula}s चा अनुक्रम}{}.
\iftag{notprvDiamond}{$\Box$ नियमाच्या आधारविधानाच्या
उजव्या बाजूला जास्तीत जास्त एकच !!{formula}~$!A$
असले पाहिजे. }{}%
\iftag{notprvBox}{$\Diamond$ नियमाच्या आधारविधानाच्या
डाव्या बाजूला जास्तीत जास्त एकच !!{formula}~$!A$
असले पाहिजे. }{}%
\iftag{prvBox}{$\Gamma$\iftag{prvDiamond}{~आणि~}{}}{}\iftag{prvDiamond}{$\Delta$}{}
रिकामे असू शकतात; अशा वेळी निष्कर्ष-क्रमवर्तीतील
त्यांचे अनुक्रमे
\iftag{prvBox}{$\Box\Gamma$\iftag{prvDiamond}{~आणि~}{}}{}\iftag{prvDiamond}{$\Diamond\Delta$}{}
हे भागही रिकामे असतात.

एका !!{formula}~$!A$ पुरतेच
\iftag{prvBox}{उजव्या बाजूला~$\Box$ जोडणे\iftag{prvDiamond}{ आणि
}{}}{}\iftag{prvDiamond}{डाव्या बाजूला~$\Diamond$ जोडणे}{}
मर्यादित ठेवणे आवश्यक आहे. जर
\iftag{prvBox}{उजवीकडे कोणत्याही संख्येने !!{formula}s ना~$\Box$
जोडण्याची\iftag{prvDiamond}{ किंवा }{}}{}\iftag{prvDiamond}{डावीकडे
कोणत्याही संख्येने !!{formula}s ना~$\Diamond$ जोडण्याची}{}
मुभा असती, तर पुढील !!{derive} करता आले असते:
  \[
    \iftag{prvBox}{
      \Axiom$!A \fCenter !A$
      \RightLabel{\RightR{\lnot}}
      \UnaryInf$\fCenter !A, \lnot !A$
      \RightLabel{$\Box*$}
      \UnaryInf$\fCenter \Box!A, \Box\lnot !A$
      \doubleLine
      \RightLabel{\RightR{\lor}}
      \UnaryInf$\fCenter \Box!A \lor \Box \lnot !A$
      \DisplayProof}{}
    \iftag{notprvBox,notprvDiamond}{}{\qquad}
    \iftag{prvDiamond}{
      \Axiom$!A \fCenter !A$
      \RightLabel{\LeftR{\lnot}}
      \UnaryInf$\lnot !A, !A \fCenter$
      \RightLabel{$\Diamond*$}
      \UnaryInf$\Diamond\lnot !A,\Diamond!A \fCenter $
      \RightLabel{\RightR{\lnot}}
      \UnaryInf$\Diamond!A \fCenter \lnot\Diamond\lnot !A$
      \RightLabel{\RightR{\lif}}
      \UnaryInf$ \fCenter \Diamond!A \lif \lnot\Diamond\lnot !A$
      \DisplayProof}{}
  \]
मात्र \iftag{prvBox}{$\Box!A \lor \Box \lnot !A$\iftag{prvDiamond}{ आणि
$\Diamond!A \lif \lnot\Diamond\lnot !A$}{}}{$\Diamond!A \lif
\lnot\Diamond\lnot !A$} यांपैकी कोणतेही सूत्र~$\Log{K}$ मध्ये
वैध नाही.

जर आधारविधानात~$!A$ व्यतिरिक्त बाजूची सूत्रेही
असू दिली, आणि \iftag{prvBox}{$\Box$ नियमाने उजवीकडे
फक्त~$!A$ ला~$\Box$ जोडू दिला\iftag{prvDiamond}{ किंवा
}{}}{}\iftag{prvDiamond}{$\Diamond$ नियमाने डावीकडे
फक्त~$!A$ ला~$\Diamond$ जोडू दिला}{} (पण बाजूच्या
सूत्रांवर कोणतीही क्रिया केली नाही), तर पुढील
!!{derive} करता आले असते:
\[
  \iftag{prvBox}{
    \Axiom$!A \fCenter !A$
    \RightLabel{\RightR{\lnot}}
    \UnaryInf$\fCenter !A, \lnot !A$
    \RightLabel{\RightR{\Exchange}}
    \UnaryInf$\fCenter \lnot !A, !A$
    \RightLabel{$\Box*$}
    \UnaryInf$\fCenter \lnot!A, \Box !A$
    \doubleLine
    \RightLabel{\RightR{\lor}}
    \UnaryInf$\fCenter \lnot!A \lor \Box !A$
    \DisplayProof}{}
  \iftag{notprvBox,notprvDiamond}{}{\qquad}
  \iftag{prvDiamond}{
    \Axiom$!A \fCenter !A$
    \RightLabel{\LeftR{\lnot}}
    \UnaryInf$\lnot !A, !A \fCenter$
    \RightLabel{$\Diamond*$}
    \UnaryInf$\Diamond\lnot !A, !A \fCenter $
    \RightLabel{\RightR{\lnot}}
    \UnaryInf$!A \fCenter \lnot\Diamond\lnot !A$
    \RightLabel{\RightR{\lif}}
    \UnaryInf$ \fCenter !A \lif \lnot\Diamond\lnot !A$
    \DisplayProof}{}
\]
मात्र \iftag{prvBox}{$\lnot!A \lor \Box !A$ (जे $!A
\lif \Box!A$ शी सममूल्य आहे)\iftag{prvDiamond}{ आणि $!A \lif \lnot\Diamond\lnot !A$
}{}}{$!A \lif \lnot\Diamond\lnot !A$} यांपैकी कोणतेही
सूत्र~$\Log{K}$ मध्ये वैध नाही.

\end{document}
